\section{Related Work}
\label{sec:related}



\paragraph{Search organization on existing graphs.}
HNSW combines hierarchical navigation and proximity-graph search~\cite{hnsw}; Faiss and Lucene implement it~\cite{faiss,lucene}. CDSM changes query allocation on existing edges. Local queues, visited state, and resumed expansion have precedents; the contribution is the specific bounded-exploration and reintegration policy.

FANNG retains unexplored edges and backtracks under a distance budget~\cite{fanng}. CDSM additionally gives backup entrances local exploration before returning unfinished work. Our control adapts FANNG search, excluding native construction and system performance.

\paragraph{Multiple paths and resumed execution.}
iQAN uses multiple paths, local queues, shared visitation, and adaptive synchronization for intra-query parallelism~\cite{iqan}. Its feedback can affect global expansion, so that principle is not unique to CDSM. Table~\ref{tab:semantics} identifies our distinction: upper-entrance exploration after a main prefix, explicit local allowances, and unfinished-state integration into one budgeted continuation. The difference concerns exploration's source, timing, allocation, and integration, beyond sequential versus parallel execution. MEP does not reproduce iQAN.

pgvector 0.8.0 retains discarded candidates and visitation to resume HNSW scans~\cite{pgvector}. CDSM first generates additional work through bounded backup exploration; resumption alone is not its novelty claim.

\paragraph{Choosing entrances.}
Oguri and Matsui select query-dependent starting points from representatives obtained by clustering~\cite{oguri}. GATE extracts hubs and learns a navigation structure using graph topology and query information~\cite{gate}. These approaches improve where exploration begins; CDSM specifies how alternative exploration and retained main-search work share a limited budget. The design dimensions are complementary, although their interaction requires evaluation. Our fixed entrance table and paid-selection extensions are not implementations of either selector, and original CDSM does not learn entrance quality.

\paragraph{Budget, workload, and quality.}
Adaptive Beam Search changes the distance-based termination rule and relates its guarantees to graph navigability~\cite{abs}. CDSM instead budgets exploratory and continuation work; its accounting limit does not imply a recall guarantee. Steiner-Hardness demonstrates the importance of graph-dependent query difficulty~\cite{steiner}, motivating evaluation beyond a workload average without making severe-failure measurement itself novel. RoarGraph addresses cross-modal workloads through query-guided index construction~\cite{roargraph}, whereas CDSM reuses existing edges. Our frozen-policy transfer results delimit its benefits: the T2I improvements do not establish a general advantage on LAION or WebVid. State interventions explain the method's behavior while its measured search-policy benefits remain the main contribution.
